\clearpage
\fancyhead[L]{\small CONNAISSANCE DES TEMS.}
\hypertarget{native-650-lower}{}
\expandafter\def\csname HistoricalNoteText37\endcsname{The present review was written in August of last year.}
\begin{center}\small Göttingische gelehrte Anzeigen. No. 6. pp. 49..56. 10 January 1828.\\[3pt]\rule{.72\linewidth}{.3pt}\end{center}

\emph{Connaissance des tems, ou des mouvemens célestes, à l'usage des astronomes et des navigateurs, pour l'an \textup{1829}. Publiée par le bureau des longitudes.} Paris 1826. Published by Bachelier. 382 pages. Large octavo.

The arrangement of the astronomical almanac remains the same as in earlier annual issues. The list of geographical longitudes and latitudes, which contains about 2000 places, may need a thorough revision or reworking. According to a remark that has become a regular feature of recent issues, the Bureau of Longitudes has commissioned a special committee to do this, and we can expect its results shortly. --- The \emph{supplements} fill 163 pages and contain the following articles. A table by \textsc{Puissant} to shorten the calculation of the altitudes of the pole and azimuths observed with the aid of the Pole Star. The astronomical observations made at many points of the trigonometric network for the new map of France prompted the calculation of these reasonably convenient auxiliary tables. They are independent of any particular declination and could therefore also be applied to off-meridian observations of other stars not too far from the pole. --- Report to the Royal Academy of Sciences by \textsc{Legendre} and \textsc{Mathieu} on a paper by Mr.\ \textsc{Puissant} concerning the determination of the figure of the Earth from geodetic and astronomical measurements. The paper deals in the first instance with the sections of two parallels measured in France and contains, according to the report, \SourceLeaf{651} methods of calculation for the problems that arise in this connection. As the paper is still to be printed in the Mémorial du dépôt de la guerre and the report does not go into detail, we cannot yet judge the value of the rules proposed by Mr.\ \textsc{Puissant}. Mr.\ \textsc{Puissant}'s method is reported only in connection with the adjustment of the angles in a system of triangles connecting two measured baselines, where the calculated value of the second baseline differs somewhat from the measured value. The method amounts to making the corrections to the angles of each triangle proportional to the error in the sum of its three angles. The authors of the report remark that this principle is convenient to apply, but that its validity has yet to be proved. We add that the \emph{correct} solution of this problem is readily found from the principles of the calculus of probabilities, which shows that the proposed principle is inadmissible. --- On the two great inequalities in the motion of Jupiter and Saturn, by \textsc{Laplace}. This paper concerns the parts of these inequalities that depend on the square of the disturbing forces. According to a remark by Mr.\ \textsc{Plana}, it is incorrect to assume the same ratio between this second part of the great equation in Saturn's motion and the second part of the great equation in the motion of Saturn as holds between the principal parts (which depend on the first power of the disturbing mass). \textsc{Laplace} works out the relation between these parts more precisely here, but the result is incompatible with the numerical values found by \textsc{Plana}. The subject thus seems to require further investigation, and it would be desirable to treat the entire theory of the mutual perturbations of Jupiter and Saturn from first principles in a way entirely different from that used so far. --- Remarks on various subjects of the Mécanique Céleste, by the same author, concern (in part as replies to some objections raised by Mr.\ \textsc{Plana}) the motion of the orbit of Saturn's outermost satellite, a periodic equation of Mercury, and the effect of the fixed stars on the solar system. --- On the measurement of an arc of the parallel at latitude $45^\circ$, by \textsc{Brousseaud} and \textsc{Nicollet}. This very interesting paper describes the operation by which the astronomical difference in longitude between Mont-Cenis and the tower of Marennes, at the mouth of the Garonne, was determined by means of gunpowder signals in 1822 and 1823. The whole arc covers 8 degrees of longitude, and the gunpowder signals were given at five intermediate points. In 1822 the difference in longitude between Mont-Cenis and Puy d'Isson in the department of Puy de Dome was determined in this way through the joint efforts of the French observers just mentioned and the Italian astronomers \textsc{Carlini} and \textsc{Plana}. In 1823 the French observers alone determined the western section, which lies wholly on French territory. The paper contains all the astronomical observations made by the French as well as the resulting values for the differences in longitude. The latter results have been extended as far as the observatory at Padua by means of similar operations carried out by the Italian astronomers in 1821 and 1824 and already made public in part. This great undertaking thus already covers an arc of 13 degrees of longitude. The arc was later extended as far as Fiume, and there is hope of extending it still farther. The geodetic determination of the arc from Marennes to Fiume rests on 106 first-order triangles, 90 of which were measured by French engineers and the rest by Austrian and Sardinian engineers. It is much to be desired that details of these triangles, too, be published soon. We have good grounds for this hope, since a notice at the end of the paper states that \emph{everything} relating to the measurement of this arc of the parallel is to be printed in the ninth volume of the Mémorial du dépôt de la guerre. The observations \SourceLeaf{652} contained in the present paper suggest a few further remarks. The determination of time was based on absolute zenith distances of stars observed with \textsc{Borda}'s repeating circles. The apparent positions of the stars were taken from \textsc{Schumacher}'s auxiliary tables. Since the utmost accuracy in determining the time is of the greatest importance in the operation under discussion, it is regrettable that portable transit telescopes were not set up at the observation stations in place of that procedure. In operations of such importance and extent, the resulting increase in cost could not have been a consideration. The times determined from stars on different sides of the meridian show differences exceeding three seconds of time, which makes simply taking the mean questionable. Moreover, at some stations the number of evenings on which the gunpowder signals were observed seems too small to inspire complete confidence in the absolute reliability of the results. Thus the difference in longitude between Solignat and La Jonchere was determined on three evenings, each time by means of 10 gunpowder signals, namely

\begin{center}\small\setlength{\tabcolsep}{3pt}
\begin{tabular}{lrll}
August & 12 & . . . . . . . . . . . . . . . . . & $6^{\mathrm m}\ 47^{\mathrm s}86$ \\
 & 19 & . . . . . . . . . . . . . . . . . & $6\quad49.81$ \\
 & $22'$ & . . . . . . . . . . . . . . . . . & $6\quad50.17$ \\
\end{tabular}\end{center}
\noindent No cause whatever could be found to explain the considerable discrepancy between the first result and the other two. The first was excluded altogether, and only the mean of the other two was used. The measurement of the arc between La Jonchere and Saint Preuil is even based solely on observations of the gunpowder signals on a single night, since the mountain Puy Cogneux, from which the signals were given, was apparently visible from St. Preuil only under stronger terrestrial refraction. Insufficient care had been taken in choosing the sites here. Since the great variability of terrestrial refraction is a well-known phenomenon that is conspicuous every day in large geodetic operations, before St. Preuil was chosen as the observing site, due consideration should have been given to the risk that the intervening high ground might obscure the signal site under weaker or moderate refraction. The small size of the staff for such extensive operations and the limited time serve as an excuse, however. Still, it remains striking that the observer stationed at St. Preuil at first lacked even \emph{precise} knowledge of the direction of the signal site, when the means of obtaining that knowledge must have been easy to find, since Puy Cogneux had been a principal triangulation point. Finally, from the entire arc from Marennes to Padua, the length of a degree of longitude along the parallel at latitude $45^\circ\,43'\,12''$ is calculated to be 77865,75 metres. By comparing the six individual parts, the reviewer finds that a mean error of 25 metres is to be expected in this result. In view of the circumstances mentioned above, however, the remaining uncertainty is likely to be greater still, so that the results for the flattening of the Earth's spheroid obtained by comparing this degree of longitude with various measured degrees of latitude still carry no great weight. — On the longitude of Port-Praslin in New Ireland, by Frigate Captain Duperrey. A new calculation of the solar eclipse observed by Veron in 1768 using \textsc{Burckhardt}'s and \textsc{Damoiseau}'s tables, and a chronometric determination made by the author aboard the Coquille, gave results in very close agreement. — On a means of eliminating the effect of capillarity in barometers, by \textsc{Laplace}. Since the mercury in the vessel is convex towards the rim, the pin that touches the surface of this mercury and whose lower end corresponds to the zero of the scale can be positioned so that the effect of capillarity is eliminated. Mr.\ \textsc{Brunard} has provided an auxiliary table for this purpose. — On the change that the constant of \SourceLeaf{653} precession undergoes as a result of the Königsberg catalogue of fundamental stars, by \textsc{Bessel}. The contents of this paper are already known from the Astronomische Nachrichten. --- Observations and elements of the comet discovered on 19 May 1825, by \textsc{Gambart}. --- On the attraction of spheroids, by \textsc{Poisson}. This substantial paper consists of two parts. The former deals with the expansion of spherical functions (as we would like to call the functions of two variable quantities which in general determine any point on the surface of a sphere) in series that proceed according to a familiar law much studied by analysts, and it is carried out with the elegance characteristic of this great geometer. An addition has been made here to the proof previously given for the fundamental theorem of this transformation, to address some proposals advanced against it. Properly understood, the addition does indeed remove all difficulties, although a fuller distinction among the infinitely small quantities involved would make the proof still more evident. At least, what a celebrated English geometer recently\HistoricalNote{37} put forward against \textsc{Poisson}'s new exposition in the May issue of the Philosophical Magazine, even if it does not touch the substance of that exposition at all, shows that the exposition could still be misunderstood. The second part contains general investigations into the attraction of arbitrary bodies, conducted with considerable analytical skill and in greater depth than before. Yet everything rests on expansion in series, and results are obtained only insofar as the body's departure from a spherical shape is assumed to be small. In the case where this condition no longer holds, in the author's words, ‘the shape of equilibrium has not yet been determined by geometers, and if this ellipsoid were the only figure in which equilibrium is possible, the strange consequence would necessarily follow that equilibrium becomes impossible at a speed of rotation which is still not the one at which the fluid begins to disperse.’ We would prefer to say that we do not yet have a mathematically rigorous proof that the ellipsoid is the only possible figure for equilibrium. For that consequence in fact contains no contradiction at all.
